% year: תשע"ח
% type: מבחן
% semester: א
% moed: ג
% number: 5א
% points: 10
% categories: antiderivatives
% source: exams/מבחנים/תשעח/מועד מיוחד תשעח.pdf

\begin{question}
{(10 נק')} חשבו את האינטגרל $\displaystyle\int\frac{x^4+4x^3+4x^2+8x+1}{(x-1)(x+2)^2}\,dx$.
\end{question}

\begin{hint}
מעלת המונה גדולה ממעלת המכנה — בצעו קודם חילוק פולינומים, ואז פרקו לשברים חלקיים מהצורה $\frac{A}{x-1}+\frac{B}{x+2}+\frac{C}{(x+2)^2}$.
\end{hint}

\begin{solution}
המכנה: $(x-1)(x+2)^2=(x-1)(x^2+4x+4)=x^3+3x^2-4$. מעלת המונה ($4$) גדולה ממעלת המכנה ($3$), לכן נחלק פולינומים:
\[ (x+1)(x^3+3x^2-4)=x^4+4x^3+3x^2-4x-4, \]
ולכן
\[ x^4+4x^3+4x^2+8x+1=(x+1)(x^3+3x^2-4)+(x^2+12x+5). \]
כלומר
\[ \frac{x^4+4x^3+4x^2+8x+1}{(x-1)(x+2)^2}=x+1+\frac{x^2+12x+5}{(x-1)(x+2)^2}. \]
פירוק לשברים חלקיים:
\[ \frac{x^2+12x+5}{(x-1)(x+2)^2}=\frac{A}{x-1}+\frac{B}{x+2}+\frac{C}{(x+2)^2}, \]
\[ x^2+12x+5=A(x+2)^2+B(x-1)(x+2)+C(x-1). \]
\begin{itemize}
  \item $x=1$: $18=9A$, כלומר $A=2$.
  \item $x=-2$: $4-24+5=-15=-3C$, כלומר $C=5$.
  \item השוואת מקדמי $x^2$: $1=A+B$, כלומר $B=-1$.
\end{itemize}
(בדיקה, מקדם חופשי: $4A-2B-C=8+2-5=5$, כנדרש.) לכן
\[ \int\frac{x^4+4x^3+4x^2+8x+1}{(x-1)(x+2)^2}\,dx=\int\left(x+1+\frac{2}{x-1}-\frac{1}{x+2}+\frac{5}{(x+2)^2}\right)dx \]
\[ \boxed{=\frac{x^2}{2}+x+2\ln|x-1|-\ln|x+2|-\frac{5}{x+2}+C} \]
\end{solution}
